\section{Datos y preprocesamiento}

\subsection{Datos sintéticos}

Las señales sintéticas se generaron integrando el sistema Balloon--Windkessel con parámetros verdaderos conocidos:

\begin{table}[H]
\centering
\caption{Parámetros utilizados para generar las señales sintéticas.}
\label{tab:synthetic_parameters}
\begin{tabular}{lc}
\toprule
Parámetro & Valor \\
\midrule
$\epsilon$ & 0,15 \\
$\kappa_s$ & 0,65 \\
$\kappa_f$ & 0,41 \\
$\tau$ & 0,98 \\
$\alpha$ & 0,32 \\
$E_0$ & 0,34 \\
$V_0$ & 0,02 \\
\bottomrule
\end{tabular}
\end{table}

Se añadió ruido estructurado compuesto por 50\% de ruido blanco, 35\% de ruido autorregresivo AR(1) y 15\% de deriva lenta. Se evaluaron niveles SNR 10, 5 y 2. El diseño completo incluyó cuatro escenarios, cinco réplicas y tres niveles de ruido, totalizando 60 señales.

La respuesta de referencia se definió operacionalmente como la respuesta del mismo sistema a un evento estandarizado de 1 s. Debido a la naturaleza no lineal del modelo, esta respuesta se utilizó como referencia comparable y no como una función impulso estricta en sentido lineal.

\subsection{Datos fMRI reales}

Se utilizaron datos de tarea motora 3T del Human Connectome Project, sujeto 100206, corridas \texttt{MOTOR\_LR} y \texttt{MOTOR\_RL}. El HCP proporciona adquisiciones multimodales y protocolos de preprocesamiento estandarizados \citep{VanEssen2013,Glasser2013}; la tarea motora forma parte del conjunto de paradigmas task-fMRI descritos en el estudio de tarea fMRI del HCP \citep{Barch2013}.

Cada corrida contiene 284 volúmenes, con TR de 0,72 s y resolución isotrópica de 2 mm. Se extrajeron señales en torno a coordenadas MNI de la corteza motora primaria izquierda $(-38,-22,56)$ y derecha $(38,-22,56)$.

Los cuatro escenarios fueron:

\begin{itemize}
    \item \texttt{MOTOR\_LR}, M1 izquierda, condición mano derecha;
    \item \texttt{MOTOR\_LR}, M1 derecha, condición mano izquierda;
    \item \texttt{MOTOR\_RL}, M1 izquierda, condición mano derecha;
    \item \texttt{MOTOR\_RL}, M1 derecha, condición mano izquierda.
\end{itemize}

\subsection{Regresión de confusores}

Para cada escenario se construyó un modelo de nuisance con las condiciones no objetivo, 12 regresores de movimiento, términos de deriva y constante. El ajuste se realizó fuera de las ventanas objetivo y posteriormente se sustrajo la contribución estimada. Dentro de cada ventana se utilizó la mediana del periodo preestímulo como línea base, evitando extrapolar una tendencia lineal que pudiera deformar la respuesta.
